\subsection{An update on moduli stabilization with antibrane uplift --- arXiv:1912.09948}

\textbf{Authors:} Emilian Dudas, Severin L\"ust

\paragraph{主な主張}
ワープ喉部の大きさを制御する軽い複素構造場が、反D3ブレーンによるuplift後に大きな超対称性破れの寄与を担い、低い超対称性破れスケールの実現には一般に大きな局在D3電荷が必要になると示した。 \cite{Dudas:2019pls}

\paragraph{新規性}
この軽い場を明示的に含め、K\"ahlerモジュライとの質量階層条件と離れたブレーン上の可視セクターへの超対称性破れの伝達を評価した。

\paragraph{位置付け}
反ブレーンupliftを含む有効理論で、積分消去の妥当性と現象論的な質量スケールを点検する研究である。
